% Complex-Analysis Toolkit  (prerequisites refresher; reference appendix for the
% single-variable complex analysis the book relies on)
\chapter{Complex-Analysis Toolkit}
\label{app:A}

This appendix collects, for reference, the results from a first course in complex
analysis that the main text uses. Proofs are in any standard text --- Ahlfors
\cite{ahlfors}, Stein--Shakarchi \cite{stein-shakarchi}, Conway \cite{conway} ---
and we state the facts precisely and note where each is invoked.

\section{Holomorphic and Meromorphic Functions}
\label{sec:A-holo}

A function \(f\) on an open \(U \subseteq \C\) is \emph{holomorphic} if it is
complex-differentiable at every point; equivalently it is given near each point by
a convergent power series. It is \emph{meromorphic} on \(U\) if it is holomorphic
off a discrete set of poles. Near a point \(p\) a meromorphic \(f \not\equiv 0\)
has a \emph{Laurent expansion}
\[
  f(z) = \sum_{k \ge m} c_k (z-p)^k, \qquad c_m \ne 0,
\]
and the integer \(m = \ord_p(f)\) is the \emph{order}: \(m > 0\) a zero, \(m < 0\)
a pole of order \(-m\), \(m = 0\) a nonzero finite value. The \emph{residue} is the
coefficient \(\Res_p f = c_{-1}\). These notions are intrinsic to a Riemann surface
because they are unchanged by holomorphic coordinate changes
(\cref{rem:div-ord-welldef}), which is what lets us speak of orders and residues on
\(T = \C/\Lambda\) throughout \crefrange{ch:elliptic_liouville}{ch:divisors}.

\section{Liouville's Theorem and the Maximum Modulus Principle}
\label{sec:A-liouville}

\begin{theorem}[Liouville]
\label{thm:A-liouville}
A bounded entire function (holomorphic on all of \(\C\)) is constant.
\end{theorem}

\begin{theorem}[Maximum modulus]
\label{thm:A-maxmod}
A nonconstant holomorphic function on a connected open set has no interior local
maximum of \(|f|\). Consequently a holomorphic function on a compact connected
Riemann surface is constant.
\end{theorem}

These are the engine of \cref{thm:ell-const}: a holomorphic elliptic function is
bounded (by periodicity and compactness of the fundamental parallelogram), hence
constant. The same statements force \(L(0) = \C\) and \(\Omega(T) = \C\,dz\) in
\cref{ch:riemann_roch}.

\section{Residues and the Residue Theorem}
\label{sec:A-residues}

\begin{theorem}[Residue theorem in the plane]
\label{thm:A-residue-plane}
If \(f\) is meromorphic on a domain containing a positively oriented simple closed
curve \(\gamma\) and its interior, with no poles on \(\gamma\), then
\[
  \frac{1}{2\pi i} \oint_\gamma f(z)\,dz = \sum_{p \text{ inside}} \Res_p f .
\]
\end{theorem}

On a compact Riemann surface there is no boundary, and the global form is a
vanishing theorem.

\begin{theorem}[Sum of residues on a compact surface]
\label{thm:A-residue-compact}
For a meromorphic \(1\)-form \(\omega\) on a compact Riemann surface \(X\),
\(\sum_{p \in X} \Res_p \omega = 0\).
\end{theorem}

The planar version drives the period-parallelogram computations of
\cref{ch:elliptic_liouville} (the boundary integral vanishes because opposite
sides cancel, leaving \(\sum \Res = 0\)); the compact version is the clean
statement behind \(\deg\divisor(f) = 0\) applied to \(df/f\) in
\cref{thm:div-degzero}. Both are global Stokes-type statements
(\cref{app:B}).

\section{The Argument Principle and Winding Numbers}
\label{sec:A-argument}

\begin{theorem}[Argument principle]
\label{thm:A-argument}
If \(f\) is meromorphic inside and on a positively oriented simple closed curve
\(\gamma\), with no zeros or poles on \(\gamma\), then
\[
  \frac{1}{2\pi i} \oint_\gamma \frac{f'(z)}{f(z)}\,dz = Z - P,
\]
the number of zeros minus the number of poles inside, with multiplicity. The
left side is the winding number of \(f \circ \gamma\) about \(0\).
\end{theorem}

Applied on a period parallelogram, this gives the valence theorem
\cref{thm:ell-valence} (\#zeros \(=\) \#poles for an elliptic function) and,
through \(\oint z\,f'/f\), the Abel relation \cref{thm:ell-abel}; the integrality
of the winding number is the \(2\pi i\,\Z\) in those proofs
(\cref{ex:ell-winding}).

\section{Pointers}
\label{sec:A-pointers}

The Weierstrass \(M\)-test (uniform convergence of a series dominated by a
convergent series of constants) justifies the convergence of the \(\wp\)-series in
\cref{prop:wp-converges}. The implicit function theorem provides holomorphic charts
on a smooth plane curve, making \(E\) a Riemann surface in \cref{sec:emb-l3o}. A
nonconstant holomorphic map between compact Riemann surfaces is a finite branched
cover and is surjective; this underlies the degree-two cover \(\wp : T \to \CP^1\)
(\cref{sec:wp-roots}) and the surjectivity of \(\Phi\) in \cref{thm:emb-iso}.
